\label[ch:goals]

As a doctoral researcher, I had the privilege to be part of the {\bf Research Center of Cosmic Rays and Radiation Events in the Atmosphere (CRREAT)}. The CRREAT programme set the following high-level aims:
i. To deepen the knowledge about the relation between atmospheric phenomena and ionising radiation.
ii. To clarify phenomena causing variations of secondary cosmic particles in the atmosphere.

In line with my expertise, the dissertation formulated the following objectives intended to contribute directly to the CRREAT aims. To connect measurements of ionising radiation with high-energy atmospheric phenomena, it was necessary to:

\nonum\sec 1. Implementation of Lightning-Detection Apparatus
Thunderstorms are known to be associated with ionising radiation, specifically {\bf terrestrial gamma-ray enhancements (TGE)} and {\bf terrestrial gamma-ray flashes (TGF)} \cite[Chilingarian_2013, Fishman_Bhat_Mallozzi_Horack_Koshut_Kouveliotou_Pendleton_Meegan_Wilson_Paciesas_et al._1994]. Lightning activity is a defining feature of thunderstorms; however, where and when exactly ionising radiation is produced within storms is not well constrained. Therefore, reliable lightning detection with appropriate spatio-temporal fidelity was a prerequisite.

At the outset I anticipated deploying a commercial solution (see Appendix \ref[app:divisek]). Field experience showed it to be insufficient, which led to a sequence of experimental approaches described in Chapter \ref[ch:trigger]. I developed a set of optical and radio techniques for lightning detection (contributing to CRREAT Aim i), implemented across {\bf mobile measurement vehicles, fixed stations, meteorological balloons, and UAVs}. Methods and field performance were reported in \cite[AMT2023] and \cite[IEEEAero2023]; the related concept of an {\bf angular-sensitive electric field mill (EFM)} that supports static and quasi-dynamic sensing is presented in \cite[JPCS2025].

\nonum\sec 2. Enhance Understanding of Spatial and Temporal Characteristics of Lightning Discharges
To characterise the lightning context of radiation events, I combined data from established detection networks \cite[Blitzortung.org, Betz2009] with the detectors and techniques developed for Objective 3. Using these resources, I conducted statistical analyses of lightning flash duration and spatio-temporal grouping.

Contrary to expectations based on literature values—average $\approx300$\,ms \cite[rakov_uman_2003] and median $\approx350$\,ms \cite[Lopez_Pineda_Montanya_Velde_Fabro_Romero_2017]—I found that the {\bf median flash duration in Central Europe is 542\,ms} (see Chapter \ref[ch:flashdur]; dataset summary in Appendix \ref[app:data]). This result has immediate implications for the design of instruments observing lightning-related transients (including ionising radiation), such as {\bf acquisition window length and trigger logic}. The dynamic observations are complemented by {\bf static/quasi-dynamic} approaches, particularly the new {\bf EFM} concept \cite[JPCS2025]. The spatio-temporal characterisation and statistics are documented in \cite[AMT2023], and similar re-assessments of flash-grouping methods have since been reported by other groups \cite[https://doi.org/10.1002/2016JD025159, 9226420].

\nonum\sec 3. Enhancement of Measurement Techniques with In-Situ Atmospheric Monitoring
To increase the probability of detecting radiation phenomena in storms and to probe high-energy processes {\it in situ}, I deployed {\bf mobile measurement vehicles}, {\bf meteorological balloons}, and {\bf UAV platforms}, and I operated {\bf remote ground-based instruments} at logistically demanding sites (mountain ridges, tall structures) lacking supporting infrastructure.

For balloon and UAV operations I designed and implemented the {\bf TF-ATMON} telemetry system and associated payloads. These platforms produced datasets enabling dosimetric analyses, including {\bf measurements of the Regener–Pfotzer maximum} and {\bf latitudinal dependence} studies, as reported in \cite[RPD2022, 10.1093/rpd/ncac299]; the telemetry architecture is summarised in \cite[AE2017].

\nonum\sec Ionising-Radiation Observations in Thunderstorms (across Objectives 1–3)
Using GEODOS and other detectors, we observed {\bf short (tens of milliseconds) enhancements} of ionising radiation that did not fit classical TGE/TGF categories, as well as {\bf radiation increases coincident with regional lightning activity yet spatially displaced by tens of kilometres} from the nearest detected discharge. These findings are reported in \cite[SilvaGabreta2024], with contemporaneous independent observations by other groups \cite[2019JD031940, 2021JD036130]. Together, they underline the necessity of multi-instrument, in-situ approaches developed and applied in this work. A broader synthesis of the methods and observations appears in CRREAT review articles \cite[Lembit1, Lembit2, Ploc1].

\nonum\sec Author’s Contribution to the Cited Publications

{\bf Monitoring of ionizing radiation from thunderstorms in Bohemian Forest using standalone device GEODOS — Silva Gabreta 30, 2024.} \cite[SilvaGabreta2024]
Co-designed GEODOS. Processed and interpreted long time series from multiple sites. (Appendix \ref[app:geodos])

{\bf Measurement of the Regener–Pfotzer maximum using different types of ionising radiation detectors and a new telemetry system TF-ATMON — Radiat. Prot. Dosim. 198(9–11): 712–719, 2022.} \cite[RPD2022]
Lead author. Approximated data for RP-maximum determination. Organised and conducted balloon flights. Conceived TF-ATMON and realised the measurement electronics. (Appendix \ref[app:rpmax])

{\bf Latitudinal effect on the position of Regener–Pfotzer maximum investigated by balloon flight HEMERA 2019 in Sweden and balloon flights FIK in Czechia — Radiat. Prot. Dosim. 199(15–16): 2041–2046, 2023.} \cite[10.1093/rpd/ncac299]
Conducted balloon flights. Built data-acquisition electronics. Evaluated telemetry datasets. (Appendix \ref[app:lat])

{\bf In situ ground-based mobile measurement of lightning events above Central Europe — Atmos. Meas. Tech. 16(2): 547–561, 2023.} \cite[AMT2023]
Lead author. Conceived the mobile measurement concept; designed/built instruments; executed field campaigns; processed long time series. (Appendix \ref[app:amt2023])

{\bf Measurements of Ionizing Radiation Generated in Thunderstorms — IEEE Aerospace Conference, 1–10, 2023.} \cite[IEEEAero2023]
Integrated mobile, balloon and GEODOS observations; contributed an imaging method using high-speed cameras and VLF antennas. (Appendix \ref[app:ieee])

{\bf Measurements with Angular Sensitive Electric Field Mill in Thunderstorms — J. Phys.: Conf. Ser. 2985(1): 012014, 2025.} \cite[JPCS2025]
Co-developed mechanical/electrical design of the angular-sensitive EFM; analysis inspired by earlier evaluation of commercial EFM-100 data. (Appendix \ref[app:jpcs])

{\bf Telemetry system for research stratospheric balloon — Int. Conf. on Applied Electronics (AE), 1–4, 2017.} \cite[AE2017]
Organised and executed the balloon flight; built supporting instruments.

{\bf Addressing challenges of high-energy atmospheric physics in Europe — Radiat. Effects Defects Solids 179(11–12): 1511–1515, 2024.} \cite[REDS2024]
CRREAT overview; contributed GEODOS, CARDOS, and mobile-vehicle measurements.

{\bf How can we simulate ionizing radiation at aviation altitudes from TGFs? — EPJ Web of Conferences 292, 09001, 2024.} \cite[EPJWeb2024]
Contributed measurements with high-speed cameras, VLF antennas, and GEODOS (simulation-focused paper).

% --- Appendix labels used:
% app:divisek, ch:trigger, ch:flashdur, app:data, app:geodos, app:rpmax, app:lat, app:amt2023, app:ieee, app:jpcs.
